\chapter{\label{cha:conclusion}Conclusions and Future Work}

\section{Conclusions}

A central finding from the uniform-grid experiments is that the poor resolution transfer of the tested FNO is mainly caused by fixed-cell spatial padding, not by the Fourier basis alone. The FNO can achieve very low error at the training resolution, but the diagnostics in Section~\ref{sec:fno_diagnostics} show that its retained-mode representation changes substantially when the same model is evaluated on a different spatial grid. In particular, the layerwise drift diagnostic in Figure~\ref{fig:fno_internal_drift} shows that fixed-cell padding introduces the first large resolution-dependent divergence from the training-resolution forward pass. The reason is that a fixed number of padded cells corresponds to a different fraction of the physical domain when the spatial resolution changes. The fixed-grid FNO experiment in Section~\ref{sec:improving_fno} supports this interpretation: mapping all input resolutions to the training resolution makes the padding operation physically consistent again and reduces \(D\) by nearly two orders of magnitude.

The CNO results show that controlling the internal representation improves resolution robustness, but also introduces interpolation and bandlimit constraints. For a convolutional encoder--decoder, a fixed internal grid prevents the physical scale of the learned kernels from changing when the input resolution changes, which explains the much flatter CNO loss curves as seen in Figures~\ref{fig:loss_vs_resolution_navier_stokes} and \ref{fig:loss_vs_resolution_darcy_flow}. However, the model then always predicts through a chosen internal resolution, so the result depends on interpolation and on the bandlimit imposed by that grid. The UND block adds further representation control by reducing nonlinear aliasing, but its empirical role is nuanced: it improves performance on Darcy flow, while CNO\_no\_UND slightly outperforms the standard CNO at high Navier--Stokes resolutions. The results therefore point to a trade-off between representation consistency and finite-grid predictive accuracy.

For non-uniform meshes, the Geo-FNO experiments show that Fourier-based operator learning can be applied directly to irregular point-cloud data. The model produced meaningful predictions on both cylinder flow and HTGR temperature data without first interpolating onto a uniform Cartesian grid. These results demonstrate applicability to irregular spatial representations, but not mesh-free discretization invariance in the strict sense: the cylinder flow point clouds change together with the geometry, while the HTGR mesh is fixed across samples.

The cylinder flow results also show that the formulation of the prediction task matters. At higher Reynolds numbers, the loss increases and the largest errors occur in the downstream wake, where vortex shedding makes the final state sensitive to the instantaneous wake phase. Geo-FNO often captures the global flow pattern and approximate wake location, but smooths the detailed oscillatory wake. This suggests that part of the remaining error is tied to one-shot prediction of an unsteady process, not only to the spatial representation used by the model.

Overall, robust neural operator surrogates require the internal representation of the model to remain consistent with the physical problem being approximated. Fixed internal resolutions are one way to improve resolution transfer, but they introduce bandlimit and interpolation assumptions. Mesh\hyp{}free models avoid regular\hyp{}grid preprocessing, but their discretization robustness is best assessed using datasets designed for controlled changes in sampling geometry. For reactor modeling, both architecture design and dataset design are therefore central to making neural operators reliable surrogate models.

\section{Future Work}

Several natural extensions of this work remain for future study.

The most immediate follow-up is to test the role of padding in the FNO more directly. The layerwise diagnostics identified fixed-cell padding as the largest source of internal representation drift under resolution transfer. A clean first test would therefore be to repeat the resolution-transfer experiments on the same datasets with padding disabled, so that the FNO is evaluated without introducing an unnecessary padded extension of the domain. For non-periodic fields, the next step would be to test physically scaled padding. A modified implementation could instead specify the padding as a fixed fraction of the domain size, so that it has the same physical meaning across resolutions. Together, these experiments would test whether the standard FNO can recover much of the robustness obtained through fixed-grid resampling while avoiding the additional interpolation step.

A second important extension is a controlled study of discretization robustness on non-uniform meshes. Unlike the uniform-grid experiments, where changing \(N\) provides a simple refinement path, non-uniform meshes require the refinement strategy itself to be specified. Future datasets should therefore contain the same physical samples represented on families of related point clouds or meshes, with controlled changes in global point count, local sampling density, or element size. For example, one could refine sampling near cylinder boundaries, wake regions, thermal boundary layers, or material interfaces while keeping the underlying physical solution fixed. This would make it possible to distinguish true robustness to discretization changes from generalization across different geometries or input parameters.

A third priority is the use of more challenging non-uniform reactor datasets. The HTGR temperature experiment demonstrates that Geo-FNO can be applied to nuclear reactor data on an irregular mesh, but the problem itself is smooth, fixed-geometry, and low-dimensional. Future work should therefore consider datasets involving stronger nonlinear coupling and more complex physical behavior, such as transient thermal-hydraulics, coupled neutronics and heat transfer or spatially varying material properties. Such datasets would provide a more meaningful test of whether neural operators can function as practical surrogate models in reactor analysis.

Another limitation of the mesh-free experiments is that only Geo-FNO was evaluated. For the uniform-grid experiments, the comparison between FNO and CNO was essential for separating the effects of global and local inductive biases. For non-uniform meshes, a similar comparison is still needed. Future work should compare Geo-FNO against locally structured mesh-free models such as CALM-PDE\footnote{Abbreviation of Continuous and Adaptive Convolutions for Latent Space Modeling of Time-dependent PDEs.} \cite{calm-pde}, as well as against the baseline of interpolating the data onto a uniform Cartesian grid before applying a standard FNO or CNO. Such comparisons would help determine whether the advantages of Geo-FNO come primarily from learning directly on the original point cloud, from the learned coordinate transform, or from the spectral structure of the architecture.

A related architectural direction is the development of a geometry-aware Convolutional Neural Operator. One possible approach would be to combine the learned coordinate transformation \(\phi^{-1}\) from Geo-FNO with the CNO architecture by first mapping the irregular point cloud to a more regular computational representation and then applying a convolutional neural operator instead of an FNO. This could preserve the local inductive bias of convolutional models while retaining the fixed internal representation and bandlimiting ideas that were effective in the uniform-grid experiments. Such an architecture would be especially relevant for cylinder flow, where the largest errors occur in localized wake structures.

The cylinder flow results also suggest that arbitrary-time or short-step prediction should be investigated for unsteady problems. The largest errors occur in cases with oscillatory wakes, where the model often captures the global wake structure but not the instantaneous wake phase. A natural extension would be to train Geo-FNO to predict short time increments and then apply it autoregressively, as is common in learned physical simulators \cite{pfaff-datasets}. This could be done using standard rollout-based training strategies, such as teacher forcing, scheduled sampling, or training on multi-step prediction losses \cite{recurrent-neural-operators}. Testing such variants would show whether the wake errors observed here are better addressed by changing the temporal formulation rather than by changing only the spatial architecture.

Overall, the central challenge remains the construction of neural operators whose internal representations remain faithful to the underlying continuous problem across changing discretizations. The results of this thesis suggest that progress on this problem will require three ingredients: architectures with controlled internal representations, datasets designed to test discretization changes cleanly, and temporal formulations that match the way the surrogate model will be used in practice.